\documentclass[10pt]{article}
\usepackage[T1]{fontenc}
\usepackage[utf8]{inputenc}
\usepackage[french]{babel}
\usepackage{amsmath,amssymb,mathrsfs,mathpazo}
\usepackage{geometry}
\geometry{a4paper,margin=2cm}
\usepackage{graphicx,array,tabularx,tabularray,colortbl,multirow,multicol,booktabs,enumitem,pifont,ulem,fancybox,diagbox,cancel}
\usepackage{pstricks,pstricks-add,pst-plot,pst-tree,pst-node,pst-text,pst-func,pst-math,pst-bspline,pst-3dplot,pst-eucl,pst-solides3d}
\let\pscalculate\undefined
\usepackage{pst-bezier}
\usepackage{tikz,pgfplots,tkz-tab}
\usetikzlibrary{arrows,arrows.meta,patterns,calc,positioning,trees,decorations.markings,shapes.geometric,shapes.misc,shapes.arrows,intersections,decorations.pathmorphing,decorations.pathreplacing,backgrounds,shadows,babel}
\pgfplotsset{compat=1.18}
\usepgfplotslibrary{fillbetween}
\usepackage[np]{numprint}
\usepackage{listings,makecell,pas-tableur,twoopt}
\let\barre\undefined
\usepackage[pstricks]{tablvar}
\definecolor{darkgreen}{RGB}{0,100,0}
\providecommand{\red}{\color{red}}
\providecommand{\blue}{\color{blue}}
\providecommand{\green}{\color{green}}
\providecommand{\black}{\color{black}}
\providecommand{\white}{\color{white}}
\providecommand{\gray}{\color{gray}}
\providecommand{\cyan}{\color{cyan}}
\providecommand{\magenta}{\color{magenta}}
\providecommand{\yellow}{\color{yellow}}
\providecommand{\violet}{\color{violet}}
\providecommand{\orange}{\color{orange}}
\providecommand{\euro}{\texteuro}
\providecommand{\EUR}{\texteuro}
\providecommand{\ieme}{\textsuperscript{e}}
\providecommand{\degres}{\ensuremath{^\circ}}
\providecommand{\degree}{\ensuremath{^\circ}}
\providecommand{\vv}[1]{\overrightarrow{#1}}
\providecommand{\Coord}[1]{#1}
\providecommand{\decosix}{}
\providecommand{\starredbullet}{\ensuremath{\bullet}}
\pagestyle{empty}
\setlength{\parindent}{0pt}
\usepackage{adjustbox}\DeclareUnicodeCharacter{2236}{:}\DeclareUnicodeCharacter{2190}{\ensuremath{\leftarrow}}
\providecommand{\R}{\mathbb{R}}
\providecommand{\N}{\mathbb{N}}
\providecommand{\Z}{\mathbb{Z}}
\providecommand{\Q}{\mathbb{Q}}
\providecommand{\D}{\mathbb{D}}
\providecommand{\C}{\mathbb{C}}
\providecommand{\e}{{\mathrm{e}}}
\providecommand{\vect}[1]{\overrightarrow{#1}}
\providecommand{\vectt}[1]{\overrightarrow{\mathrm{#1}}}
\providecommand{\vvt}[1]{\overrightarrow{\mathrm{#1}}}
\providecommand{\np}[2][]{{#2~#1}}
\providecommand{\vv}[1]{{\overrightarrow{#1}}}
\renewcommand{\d}{{\mathrm{d}}}
\providecommand{\degres}{{^\circ}}
\providecommand{\degre}{{^\circ}}
\providecommand{\blue}{\color{blue}}
\providecommand{\red}{\color{red}}
\providecommand{\enskip}{\quad}
\providecommand{\llbracket}{\lbrack\!\lbrack}
\providecommand{\rrbracket}{\rbrack\!\rbrack}
\providecommand{\euro}{{\text{€}}}
\providecommand{\no}{{\text{n°}}}
\providecommand{\interval}[2]{{[#1;#2]}}
\providecommand{\intervFO}[2]{{[#1;#2[}}
\providecommand{\dotfill}{\ldots}

\providecommand{\ds}{\displaystyle}
\providecommand{\textding}[1]{#1}
\providecommand{\Oij}{\(\mathrm{O};\vec{\imath},\vec{\jmath}\)}
\providecommand{\Oijk}{\(\mathrm{O};\vec{\imath},\vec{\jmath},\vec{k}\)}
\providecommand{\Ouv}{\(\mathrm{O};\vec{u},\vec{v}\)}
\providecommand{\intervFF}[2]{\ensuremath{[#1;#2]}}
\renewcommand{\euro}{\ifmmode\text{\texteuro}\else\texteuro\fi}
\DeclareUnicodeCharacter{00B0}{\ifmmode^{\circ}\else\textdegree\fi}
\DeclareUnicodeCharacter{20AC}{\ifmmode\text{\texteuro}\else\texteuro\fi}
% Énoncés : Éducation nationale. Transcription et corrigés : APMEP, Denis Vergès et contributeurs. Découpage et classement : CoopMaths, Sébastien Lozano, Éric Elter et contributeurs. Adaptation de la présentation : Maths Explorer.
% CC BY-SA : https://coopmaths.fr/www/about/
% https://coopmaths.fr/alea/static/bac/2024/tex/bac_2024_06_sujet2_polynesie_1.tex
\begin{document}
\noindent\textit{APMEP / CoopMaths. CC BY-SA. Présentation adaptée par Maths Explorer.}\par\medskip
\medskip

Un sondage réalisé en France fournit les informations suivantes :

\begin{itemize}[label=$\bullet~$]
\item 60\,\% des plus de 15 ans ont l'intention de regarder les jeux Olympiques et Paralympiques (JOP) de Paris 2024 à la télévision ;
\item parmi ceux qui ont l'intention de regarder les JOP, 8 personnes sur 9 déclarent pratiquer une activité sportive régulière.
\end{itemize}

On choisit au hasard une personne de plus de 15 ans. On considère les évènements suivants :

\begin{itemize}[label=$\bullet~$]
\item $J$ : \og la personne a l'intention de regarder les JOP Paris 2024 à la télévision \fg\ ;
\item $S$ : \og la personne choisie déclare pratiquer une activité sportive régulière \fg.
\end{itemize}

On note $\overline{J}$ et $\overline{S}$ leurs évènements contraires.

\smallskip

\emph{Dans les questions $1$. et $2$., les probabilités seront données sous la forme d'une fraction irréductible.}

\medskip

\begin{enumerate}
	\item Démontrer que la probabilité que la personne choisie ait l'intention de regarder les JOP de Paris 2024 à la télévision et déclare pratiquer une activité sportive régulière est de $\dfrac{8}{15}$.
	
\emph{On pourra s'appuyer sur un arbre pondéré.}
\end{enumerate}

Selon ce sondage, deux personnes sur trois parmi les plus de 15 ans déclarent pratiquer une activité sportive régulière.

\begin{enumerate}[resume]
\item 
	\begin{enumerate}
		\item Calculer la probabilité que la personne choisie n'ait pas l'intention de regarder les JOP de Paris 2024 à la télévision et déclare pratiquer une activité sportive régulière.
		\item En déduire la probabilité de $S$ sachant $\overline{J}$ notée $P_{\overline{J}}(S)$.
	\end{enumerate}
\end{enumerate}

\medskip

\emph{Dans la suite de l'exercice, les résultats seront arrondis au millième.}

\begin{enumerate}[resume]
\item Dans le cadre d'une opération de promotion, $30$ personnes de plus de $15$~ans sont choisies au hasard.
	
On assimile ce choix à un tirage avec remise.
	
On note $X$ la variable aléatoire qui donne le nombre de personnes déclarant pratiquer une activité sportive régulière parmi les $30$ personnes.
	\begin{enumerate}
		\item Déterminer la nature et les paramètres de la loi de probabilité suivie par $X$.
		\item Calculer la probabilité qu'exactement $16$~personnes déclarent pratiquer une activité sportive régulière parmi les $30$~personnes.
		\item La fédération française de judo souhaite offrir une place pour la finale de l'épreuve par équipe mixte de judo à l'Arena Champ-de-Mars pour chaque personne déclarant pratiquer une activité sportive régulière parmi ces $30$~personnes.
		
Le prix d'une place s'élève à 380~\euro{} et on dispose d'un budget de \np{10000}~euros pour cette opération.
		
Quelle est la probabilité que ce budget soit insuffisant ?
	\end{enumerate}
\end{enumerate}

\end{document}
